\chapter{Конструкторская часть}
% #todo тут рассмотрены...

\section{Схемы алгоритмов}
Ниже приведены схемы алгоритмов отрисовки сцены и генерации городской среды.
\begin{figure}[H]
	\centering
	\includegraphics[width=0.6\textwidth]{cityGeneration}
	\caption{Схема алгоритма создания сцены городской среды}
	\label{fig:cityGenerationAlgScheme}
\end{figure}
\begin{figure}[H]
	\centering
	\includegraphics[width=0.9\textwidth]{RendererAlgo}
	\caption{Схема алгоритма отрисовки объектов}
	\label{fig:RendererAlgo}
\end{figure}

Распараллеливание вычислений производится при помощи обработки пикселей при растеризации в различных потоках.

\section{Используемые структуры данных}

Для реализации программного обеспечения были использованы следующие структуры данных.

\subsection*{Камера}

Камера описывает точку зрения, с которой производится отрисовка сцены. Структура содержит следующие поля:
\begin{itemize}
	\item трехмерный вектор позиции камеры в пространстве;
	\item целевая точка, на которую направлена камера;
	\item вектор вверх.
\end{itemize}

Для построения карт теней используется ортографическая камера, которая обеспечивает параллельное проецирование без перспективных искажений, что критически важно для корректного вычисления теней.

\subsection*{Графический объект}

Графический объект представляет собой базовую структуру для всех объектов, которые могут быть отрисованы в сцене. Структура содержит:
\begin{itemize}
	\item массив точек -- список вершин, составляющих геометрию объекта;
	\item массив граней -- список граней, определяющих полигональную сетку объекта.
\end{itemize}

\subsection*{Грань}

Грань описывает грань полигональной сетки и содержит:
\begin{itemize}
	\item три индекса вершин, образующих треугольную грань;
	\item нормальный вектор грани, используемый для расчёта освещения;
	\item цвет грани.
\end{itemize}

\subsection*{Сцена}

Сцена представляет собой контейнер для всех объектов, находящихся в виртуальном трехмерном пространстве. Структура включает:
\begin{itemize}
	\item массив графических объектов;
	\item ссылку на активную камеру;
	\item массив источников света.
\end{itemize}


\section*{Вывод}
В конструкторской части представлены схемы алгоритмов генерации городской среды и отрисовки объектов. Описаны ключевые структуры данных, используемые в программной реализации: камера, графический объект, грань и сцена. Определен подход к распараллеливанию вычислений путём обработки пикселей в различных потоках при растеризации.